%&program=xetex
%&encoding=UTF-8 Unicode

\TeXXeTstate=1
\nopagenumbers \frenchspacing
%\font\title="Nafees Web Naskh pakistani BTK:script=arab" at 72pt
%\font\heading="Nafees Tahreer Naskh v1.0:script=arab" at 18pt
\font\body="Nafees Tahreer Naskh v1.0:script=arab" at 15pt

\font\title="Nafees Web Naskh pakistani BTK:script=arab" at 72pt
\font\heading="Nafees Web Naskh pakistani BTK:script=arab" at 18pt
\font\bodybold="Nafees Web Naskh pakistani BTK:script=arab" at 15pt
\parindent=0.5in \baselineskip=22pt \lineskiplimit=-1000pt

\def\ttl#1{\beginR\title #1\endR}
\def\hdr#1{\beginR\heading #1\endR}
\def\pb#1{\beginR\body #1\endR}


\centerline{\hdr{ابنِ صفی کا تعارف}}
\everypar={\setbox0=\lastbox \beginR \box0 }
\bigskip\bigskip\bigskip

\beginR
\body

ابنِ صفی ہندستان کے شہر اِلہ آباد(اُتر پردیش) میں۲۶ جولائی۱۹۲۴ء کو پیدا ہوئے۔ اُنکا اصل نام اِسرار احمد تھا۔ بی اے کی ڈگری آگرہ یونیورسٹی سے حاصل کی۔ اُنہوں نے لکھنے کا آغاز ۱۹۴۰ء میں افسانے اور طنزومزاح لکھنے سے کیا۔ پھر اُنہوں نے ۱۹۵۰ء میں ناول نگاری شروع کی، اُس وقت وہ سیکنڈری سکول ٹیچر کی حیثیت سے کام کر رہے تھے۔ آزادی سے پہلے اور بعد میں اپنی حریت پسند طبیعت کے باعث حکومت کی نظروں میں آ چکے تھے، اسلیے ۱۹۵۲ء میں آپ کراچی منتقل ہو گئے۔

ابنِ صفی نے جاسوسی ناول نگاری کسی صاحب کے دعوے کے بعد کی تھی کہ برِصغیر میں جاسوسی ناول فحش نگاری کی وجہ سے مقبول ہیں۔ انہوں نے اسے غلط ثابت کرنے کیلئے  پہلے \bodybold جاسوسی دنیا \body اور پھر \bodybold عمران سیریز \body  شروع کی۔ ابنِ صفی کے بقول، اُنکے صرف آٹھ ناولوں کے مرکزی خیال کسی سے مُستعار لیے ہیں، باقی کے ۲۴۵ ناول مکمل طور پر انکے اپنے ہیں۔

۱۹۶۰ء سے ۱۹۶۳ء تک آپ Schziophrenia کے مریض رہے، لیکن ۱۹۶۳ء میں ہی اس مرض سے نہ صرف سنبھالا لیا بلکہ عمران سیریز کے بہترین ناول \bodybold ڈیڑہ متوالے \body بھی لکھا۔ ۱۹۷۰ء میں پاکستانی اِنٹیلی جنس ISI کو جاسوسی کے حوالے سے غیر رسمی مشاورت بھی دی۔

۲۶ جولائی ۱۹۸۰ء میں (یعنی اپنی سالگرہ کے دن) وہ اس جہان سے رخصت ہوئے، وجہِ انتقال موضی مرض کینسر تھا۔ ابنِ صفی کا لگایا ہوا پودا عمران سیریز اسقدر تناور ہوا کہ آج تک یہ ثمر بار ہے اور کئی ایک مصنفین اس سلسلے کو آگے بڑھا رہے ہیں۔

\endR

\break

\vfil
\centerline{\ttl{خوفناک عمارت}}
\bigskip\bigskip\bigskip

\centerline{\hdr{عمران سیریز کا پہلا ناول}}
\break


\beginR
\body

"ناپ کر دیکھ لو میری جان، اگر غلط نکلے تو میرا قلم سر کر دینا۔۔۔آں۔۔۔شائد میں غلط بول گیا۔۔۔میرےقلم پہ سر رکھ دینا۔۔۔" عمران نے کہا اور اِدھر اُدھر دیکھنے لگا۔ اس نے ایک طرف پڑا ہوا تنکا اُٹھایا اور پھر جھک کر زخموں کا درمیانی فاصلہ ناپنے لگا۔ فیاض اسے حیرت سے دیکھ رہا تھا۔

"لو" عمران اسے تنکا پکڑاتے ہوئے بولا۔ "اگر یہ تنکا پانچ انچ کا نہ نکلے تو کسی کی داڑھی میں تلاش کر لینا۔"

"مگر اس کا مطلب!" فیاض کچھ سوچتا ہوا بولا۔

"اس کا مطلب یہ کہ قاتل و مقتول دراصل عاشق و معشوق تھے۔"

"عمران پیارے ذرا سنجیدگی سے۔"

"یہ تنکا بتاتا ہے کہ یہی بات ہے۔" عمران نے کہا "اور اُردو کے پرانے شعراء کا بھی یہی خیال ہے۔ کسی کا بھی دیوان اُٹھا کر دیکھ لو! دو چار شعر اس قسم لے ضرور مل جائیں گے جن سے میرے خیال کی تائید ہو جائے گی۔ چلو ایک شعر سن ہی لو۔"

موچ آئے نہ کلائی میں کہیں	\hfil سخت جاں ہم بھی بہت پیارے

"مت بکواس کرو۔ اگر میری مدد نہیں کرنا چاہتے تو صاف صاف کہ دو۔" فیاض بگڑ کر بولا۔

"فاصلہ تم نے ناپ لیا! اب تم ہی بتاؤ کہ کیا بات ہو سکتی ہے" عمران نے کہا۔

فیاض کچھ نہ بولا۔

"ذرا سوچو تو۔" عمران پھر بولا۔ "ایک عاشق ہی اُردو شاعری کے مطابق اپنے محبوب کو اِس بات کی اجازت دےسکتا ہے کہ جسطرح چاہے اسے قتل کرے۔ قیمہ بنا کر رکھ دے یا ناپ ناپ کر سلیقے سے زخم لگائے، یہ زخم بدحواسی کا نتیجہ نہیں۔ لاش کی حالت بھی یہ نہیں بتاتی کہ مرنے سے پہلے مقتول کو کسی سے جدوجہد کرنی پڑی ہو۔ بس ایسا معلوم ہوتا ہے جیسے چپ چاپ لیٹ کر اس نے کہاجو مزاجِ یار میں آئے۔۔۔"

"پرانی شاعری اور حقیقت میں کیا لگاؤ ہے؟" فیاض نے پوچھا۔

"پتہ نہیں۔" عمران پر خیال انداز میں سر ہلا کر بولا۔ "ویسے اب تم پوری غزل سنا سکتے ہو۔ مقطع میں عرض کر دوں گا۔"

فیاض تھوڑی دیر خاموش رہا پھر بولا۔ "یہ عمارت تقریباً پانچ سال سے خالی رہی ہے! ۔۔۔ ویسے ہر جمعرات کو صرف چند گھنٹوں کے لئے اسے کھولا جاتا ہے۔"

"کیوں؟"

"یہاں دراصل ایک قبر ہے جس کے متعلق مشہور ہے کہ وہ کسی شہید کی ہے، چنانچہ ہر جمعرات کو ایک شخص اسے کھول کر اس قبر کی جاروب کشی کرتا ہے۔"

"چڑھاوے وغیرہ چڑھتے ہوں گے۔" عمران نے پوچھا۔

"نہیں ایسی کوئی بات نہیں۔ جن لوگوں کا یہ مکان ہے وہ شہر میں رہتے ہیں اور ان سے میرے قریبی تعلقات ہیں۔ انہوں نے ایک آدمی کو اسی لئے رکھ چھوڑا ہے کہ وہ ہر جمعرات کو قبر کی دیکھ بال کر لیا کرے! ۔۔۔ یہاں معتقدین کی بھیڑ نہیں ہوتی۔ بہرحال آدمی آج دوپہر کو جب یہاں آیا تو اس نے یہ لاش دیکھی۔"

"تالا بند تھا؟" عمران نے پوچھا۔




\endR

\bye